\begin{frame}[shrink]
	\frametitle{Cada codificação tem seu próprio piso de erro}
	\small
	\par Mesmo antes do treino, converter um valor contínuo em pulsos pode perder precisão. Esse \textbf{piso de codificação} não é ruído do sensor nem erro do classificador: é o erro que a representação e sua leitura já permitem.
	\par Exemplo idealizado: $x$ uniforme em $[0,1]$, $T=16$ passos e a mesma métrica (erro RMS de reconstrução):
	\begin{center}
		\begin{tabular}{lcc}
			\toprule
			\textbf{Codificação} & \textbf{De onde vem o erro} & \textbf{Piso RMS} \\
			\midrule
			Direta & $x$ entra sem quantização em pulsos & $0$ \\
			Latência & arredondamento para um dos 16 instantes & $\frac{1}{15\sqrt{12}}\approx 0{,}019$ \\
			Poisson & contagem aleatória de pulsos & $\sqrt{\frac{1}{6\cdot16}}\approx 0{,}102$ \\
			\bottomrule
		\end{tabular}
	\end{center}
	\begin{itemize}
		\item O zero da codificação direta vale só para esta etapa ideal de entrada e reconstrução; não significa que a rede inteira terá erro zero.
		\item \textbf{Por que importa?} Se um modelo começa com um piso maior, comparar apenas seu erro final pode atribuir à arquitetura uma diferença que veio da codificação.
		\item Compare as redes com o mesmo $T$, a mesma métrica e em relação ao piso de cada codificação. A escolha entre códigos também envolve latência, energia e robustez, não só erro \cite{eshraghian2023training, liu2023firstspike}.
	\end{itemize}
\end{frame}
